\sect{भाद्रपद-कृष्णाजा-एकादशी-माहात्म्यम्}
\label{sec:padma-bhadrapada-krishnaja}

\ganapatyadiDhyanam
\padmaPuranam

\dnsub{अथ अष्टपञ्चाशत्तमोऽध्यायः॥५८}

\uvacha{युधिष्ठिर उवाच}

\twolineshloka
{भाद्रस्य कृष्णपक्षे तु किन्नामैकादशीभवेत्}
{एतदिच्छाम्यहं श्रोतुं कथयस्व जनार्दन}% ॥१॥

\uvacha{श्रीकृष्ण उवाच}

\twolineshloka
{शृणुष्वैकमना राजन् कथयिष्यामि विस्तरात्}
{अजेति नामतः प्रोक्ता सर्वपापप्रणाशिनी}% ॥२॥

\twolineshloka
{पूजयित्वा हृषीकेशं व्रतमस्यां करोति यः}
{पापानि तस्य नश्यन्ति व्रतस्य श्रवणादपि}% ॥३॥

\twolineshloka
{नातः परतरा राजन् लोकद्वयहिताय वै}
{सत्यमुक्तं मया ह्येतन्नासत्यं मम भाषितम्}% ॥४॥

\twolineshloka
{हरिश्चन्द्र इति ख्यातो बभूव नृपतिः पुरा}
{चक्रवर्ती सत्यसन्धः समस्ताया भुवः पतिः}% ॥५॥

\twolineshloka
{कस्यापि कर्मणः प्राप्तौ राज्यभ्रष्टो बभूव सः}
{विक्रीतौ वनितापुत्रौ स चकारात्मविक्रयम्}% ॥६॥

\twolineshloka
{पुल्कसस्य च दासत्वं गतो राजा स पुण्यकृत्}
{सत्यमालम्ब्य राजेन्द्र मृतचैलापहारकः}% ॥७॥

\twolineshloka
{सोऽभवन्नृपतिश्रेष्ठो न सत्याच्चलितस्तथा}
{एवं च तस्य नृपतेर्बहवो वत्सरा गताः}% ॥८॥

\twolineshloka
{ततश्चिन्तापरो राजा स बभूवातिदुःखितः}
{किं करोमि क्व गच्छामि निष्कृतिर्मे कथं भवेत्}% ॥९॥

\twolineshloka
{इति चिन्तयतस्तस्य मग्नस्य वृजिनार्णवे}
{आजगाम मुनिः कश्चिज्ज्ञात्वा राजानमातुरम्}% ॥१०॥

\twolineshloka
{परोपकारणार्थाय निर्मिता ब्रह्मणा द्विजाः}
{स तं दृष्ट्वा द्विजवरं ननाम नृपसत्तमः}% ॥११॥

\twolineshloka
{कृताञ्जलिपुटो भूत्वा गौतमस्याग्रतः स्थितः}
{कथयामास वृत्तान्तमात्मनो दुःखसंयुतम्}% ॥१२॥

\twolineshloka
{श्रुत्वा नृपतिवाक्यानि गौतमो विस्मयान्वितः}
{उपदेशं नृपतये व्रतस्यास्य ददौ मुनिः}% ॥१३॥

\twolineshloka
{मासि भाद्रपदे राजन् कृष्णपक्षेति शोभना}
{एकादशी समायाता  अजा नामेति पुण्यदा}% ॥१४॥

\twolineshloka
{अस्याः कुरु व्रतं राजन् पापस्यान्तो भविष्यति}
{तव भाग्यवशादेषा सप्तमेऽह्नि समागता}% ॥१५॥

\twolineshloka
{उपवासपरो भूत्वा रात्रौ जागरणं कुरु}
{एवमस्या व्रते चीर्णे तव पापक्षयो ध्रुवम्}% ॥१६॥

\twolineshloka
{तव पुण्यप्रभावेण चागतोऽहं नृपोत्तम}
{इत्येवं कथयित्वा च मुनिरन्तरधीयत}% ॥१७॥

\twolineshloka
{मुनिवाक्यं नृपः श्रुत्वा चकार व्रतमुत्तमम्}
{कृते तस्मिन् व्रते राज्ञः पापस्यान्तोऽभवत् क्षणात्}% ॥१८॥

\twolineshloka
{श्रूयतां राजशार्दूल प्रभावोऽस्य व्रतस्य च}
{यद् दुःखम्बहुभिर्वर्षैर्भोक्तव्यन्तत्क्षयोभवेत्}% ॥१९॥

\twolineshloka
{निस्तीर्णदुःखो राजासीद् व्रतस्यास्य प्रभावतः}
{पत्न्या सह समायोगं पुत्रजीवनमाप सः}% ॥२०॥

\twolineshloka
{दिवि दुन्दुभयो नेदुः पुष्पवर्षमभूद्दिवः}
{एकादश्याः प्रभावेन प्राप्यराज्यमकण्टकम्}% ॥२१॥

\twolineshloka
{स्वर्गं लेभे हरिश्चन्द्रः सपुरः सपरिच्छदः}
{ईदृग्विधं व्रतंराजन् ये कुर्वन्ति च मानवाः}% ॥२२॥

\twolineshloka
{सर्वपापविनिर्मुक्तास्त्रिदिवं यान्ति ते नृप}
{पठनाच्छ्रवणाद्वाऽपि अश्वमेधफलं लभेत्}% ॥२३॥

॥इति श्रीपाद्मे महापुराणे पञ्चपञ्चाशत्साहस्र्यां संहितायामुत्तरखण्डे उमापतिनारदसंवादान्तर्गतकृष्णयुधिष्ठिरसंवादे भाद्रपद-कृष्णाजा-एकादशी-माहात्म्यं नाम अष्टपञ्चाशत्तमोऽध्यायः॥५८॥

\genMangalaShloka

\hyperref[sec:ekadashi_mahatmyam_padma_puranam]{\closesub}
\clearpage

